\chapter{ระบบเชื่อมโยงข้อมูลคลาวด์และ Google Drive}
\label{chap:cloud_sync}

\section{สถาปัตยกรรมการแลกเปลี่ยนข้อมูลแบบเปิด (Open Data Pipeline)}
เพื่อให้ข้อมูลคุณภาพดินที่ตรวจวัดได้จากภาคสนามสามารถนำไปใช้ประโยชน์ต่อในงานวิเคราะห์เชิงสถิติขั้นสูง การสร้างฐานข้อมูลสารสนเทศภูมิศาสตร์ (GIS) หรือการประมวลผลข้อมูลขนาดใหญ่ (Big Data Analytics) ระบบ SoilpHTxAI จึงได้รับการออกแบบระบบส่งออกข้อมูล (Data Export Pipeline) ที่มีความยืดหยุ่นและรองรับรูปแบบมาตรฐานสากล

ระบบผสานการทำงานระหว่างหน่วยความจำภายในเครื่อง (Local File Cache) และระบบแชร์ของระบบปฏิบัติการ (Native Android Share Sheet) ทำให้สามารถส่งต่อชุดข้อมูลไปยัง Google Drive, อีเมลของสถาบัน, แอปพลิเคชันส่งข้อความ หรือเซิร์ฟเวอร์ฐานข้อมูลกลางได้อย่างไร้รอยต่อ

\begin{figure}[H]
\centering
\resizebox{0.88\textwidth}{!}{
\begin{tikzpicture}[
    node distance=1.5cm,
    step_node/.style={rectangle, draw=rbruNavy, fill=white, rounded corners=6pt, text width=3.2cm, text centered, minimum height=1.2cm, line width=1pt, drop shadow={opacity=0.15}},
    cloud_node/.style={rectangle, draw=rbruBlue, fill=rbruBlue!10, rounded corners=6pt, text width=3.2cm, text centered, minimum height=1.2cm, line width=1.5pt, drop shadow={opacity=0.2}},
    arrow/.style={draw=rbruNavy, -{Latex[length=3mm]}, line width=1.5pt}
]
    \node [step_node] (data) {\textbf{ข้อมูลในแอป}\\Calibration / Field};
    \node [step_node, right=1.6cm of data] (gen) {\textbf{จัดทำโครงสร้าง}\\ไฟล์ .CSV / .JSON};
    \node [step_node, right=1.6cm of gen] (share) {\textbf{Android Share}\\Sheet Service};
    \node [cloud_node, below=1.2cm of share] (gdrive) {\textbf{Google Drive}\\จัดเก็บบนคลาวด์};
    \node [cloud_node, left=1.6cm of gdrive] (server) {\textbf{Cloud REST API}\\ฐานข้อมูลกลาง};

    \path [arrow] (data) -- (gen);
    \path [arrow] (gen) -- (share);
    \path [arrow] (share) -- (gdrive);
    \path [arrow] (share) -| (server);
\end{tikzpicture}
}
\caption{สถาปัตยกรรมท่อส่งข้อมูลการสำรวจดินขึ้นสู่ Google Drive และคลาวด์}
\label{fig:cloud_pipeline}
\end{figure}

\section{รูปแบบโครงสร้างข้อมูลส่งออก}

\subsection{โครงสร้างไฟล์ข้อมูลสอบเทียบในห้องแล็บ (.CSV)}
ไฟล์ข้อมูลสอบเทียบถูกส่งออกเป็นมาตรฐาน Comma-Separated Values ประกอบด้วยคอลัมน์สำคัญ ได้แก่ รหัสตัวอย่าง (\texttt{id}), ศักย์ไฟฟ้า (\texttt{potential\_mv}), อุณหภูมิ (\texttt{temperature\_c}), ค่ามาตรฐานบัฟเฟอร์ (\texttt{standard\_ph}), ค่าสูตรเนิร์นสต์ (\texttt{nernst\_ph}), ค่าทำนายของ AI (\texttt{ai\_predicted\_ph}), ค่าความคลาดเคลื่อนดั้งเดิม (\texttt{traditional\_error}), ค่าความคลาดเคลื่อน AI (\texttt{ai\_error}), การแบ่งส่วนชุดข้อมูล (\texttt{split}), และเวลาที่บันทึก (\texttt{timestamp})

\subsection{โครงสร้างไฟล์ข้อมูลแปลงสำรวจดินภาคสนาม (.JSON)}
ข้อมูลแปลงดินสำรวจ 30 แปลงถูกส่งออกในรูปแบบ JavaScript Object Notation ที่มีโครงสร้างลำดับชั้น (Hierarchical JSON) เพื่อความสะดวกในการนำเข้าสู่ระบบแผนที่ GIS และการประมวลผลต่อ ดังแสดงในตัวอย่างรหัสที่ \ref{lst:json_sample}

\begin{codebox}[language=Java]{ตัวอย่างโครงสร้างข้อมูลแปลงดินภาคสนามรูปแบบ JSON}
{
  "project": "SoilpHTxAI - Rambhai Barni Rajabhat University",
  "exported_at": "2026-10-01T05:30:00.000Z",
  "record_count": 30,
  "records": [
    {
      "sample_id": "FLD-0001",
      "timestamp": "2026-10-01T04:15:22.000Z",
      "potential_mv": 128.4,
      "temperature_c": 28.5,
      "raw_ph": 4.85,
      "nernst_ph": 4.83,
      "ai_ph": 4.88,
      "site_name": "Nong-Or Durian Orchard Block A",
      "province": "Chanthaburi",
      "latitude": 12.671920,
      "longitude": 102.193240,
      "soil_type": "Sandy Loam",
      "durian_status": "Severe Acidic (Critical)",
      "recommended_lime_kg_per_rai": 250.0
    }
  ]
}
\end{codebox}
\label{lst:json_sample}

\clearpage
